\begin{bt}%[1D1V2-5]%[ma:b532365c]
% nguon: Archive/SP_KhoLatex/ID6_old/1K1-ot.tex #3778
%[Dự án 2 - TLDH - TeamTeXHoa - Quang Ngô]
    Khi một tia sáng truyền từ ông khí vào mặt nước thì một phần tia sáng bị phản xạ trên bề mặt, phần còn lại bị khúc xạ như trong Hình. Góc tới $i$ liên hệ với góc khúc xạ $r$ bởi Định luật khúc xạ ánh $\dfrac{\sin i}{\sin r}=\dfrac{n_2}{n_1}. $
    \begin{center}
        \begin{tikzpicture}[line join = round, line cap = round, >=stealth, font=\footnotesize, scale=1]
            
            \path 
            (-1,0) coordinate (A)
            (5,0) coordinate (B)
            (2,0) coordinate (C)
            (-1,2) coordinate (D)
            (5,2) coordinate (E)
            (2,3) coordinate (F)
            (2,-3) coordinate (G)
            (4,-3) coordinate (H)
            ;
            \fill[color=blue!20] (A)--(B)--(5,-3.2)--(-1,-3.2)--(A);
            \draw
            (A)--(B)
            ($(C)!.7!(H)$)--(H)node[right]{$R$}
            ($(D)!.3!(C)$)--(C)node[below left]{$I$}
            ($(C)!.7!(E)$)--(E)node[above]{$S'$}
            pic[draw,black,angle radius=0.3cm] {angle = F--C--D}
            (C)++(120:.5)node{$i$}
            pic[draw,black,angle radius=0.4cm] {angle = E--C--F}
            (C)++(60:.6)node{$i'$}
            pic[draw,black,angle radius=0.3cm] {angle = G--C--H}
            pic[draw,black,angle radius=0.4cm] {angle = G--C--H}
            (C)++(-70:.6)node{$r$}
            ;
            \draw[->]
            (C)--($(C)!.7!(H)$);
            \draw[->]
            (D)node[above]{$S$}--($(D)!.3!(C)$);
            \draw[->]           
            (C)--($(C)!.7!(E)$);
            \draw[dashed] (F)node[left]{$N$}--(G)node[left]{$N'$};
            \draw 
            (A)++(60:.5) circle (2mm)node{$1$}
            (A)++(-60:.5) circle (2mm)node{$2$}
            ;
        \end{tikzpicture}
    \end{center}
    Ở đây, $n_1$ và $n_2$ tương ứng là chiết suất của môi trường 1 (không khí) và môi trường 2 (nước). Cho biết góc tới $i=50^\circ$, hãy tính góc khúc xạ, biết rằng chiết suất của không khí bằng 1 còn chiết suất của nước là $1{,}33$.
    \loigiai{
        Theo bài ra ta có: $ i=50^\circ$, $n_1=1$, $n_2=1{,}33$, thay vào $\dfrac{\sin i}{\sin r}=\dfrac{n_2}{n_1}$ ta được:\\
        $\dfrac{\sin{50}^\circ}{\sin r}=\dfrac{1{,}33}{1}$ (đk $\sin r\neq 0$)\\
        $\Rightarrow \sin r=\dfrac{\sin50^\circ}{1{,}33}$$\Leftrightarrow \sin r\approx 0{,}57597$ (thỏa mãn điều kiện)\\
        $\Leftrightarrow \sin r\approx \sin\left(35^\circ10^\prime\right)$\\
        $\Leftrightarrow \hoac{
            &r\approx 35^\circ10^\prime+k360^\circ \\
            &r\approx 180^\circ-35^\circ10^\prime+k360^\circ 
        }\left(k\in \mathbb{Z}\right) \\
        \Leftrightarrow \hoac{
            &r\approx 35^\circ10^\prime+k360^\circ \\
            &r\approx 144^\circ50^\prime+k360^\circ 
        }\left(k\in \mathbb{Z}\right)$ \\
        Mà $0^\circ<r<90^\circ$ nên $r\approx 35^\circ10^\prime$.\\
        Vậy góc khúc xạ $r\approx 35^\circ10^\prime.$
    }
\end{bt}

\begin{bt}%[1D3V2-5]%[ma:3666bb2d]
% nguon: Archive/SP_SauLuoi_2024-2025/SP_HKI/1-TL-TN-DS-THPT-Hoang-Viet-Daklak-HKI-NH24-25.tex #48
%[Dự án đề kiểm tra Toán 11 HKI NH24-25-Dot5-Thanh Phong]%[TH,THCS & THPT Hoàng Việt - Đăklăk]
Tính giới hạn $\underset{x\to-\infty}{\lim}\,\left(\sqrt{4x^2-x+2}+2x-1\right)$ .
\loigiai{
Ta có
	\begin{eqnarray*}
	& &\lim\limits_{x\to-\infty}\left(\sqrt{4x^2-x+2}+2x-1\right)\\
	&=&\lim\limits_{x\to-\infty}\left(\dfrac{4x^2-x+2-\left(2x-1\right)^2}{\sqrt{4x^2-x+2}-2x+1}\right)\\
	&=&\lim\limits_{x\to-\infty}\left(\dfrac{3x+1}{\sqrt{4x^2-x+2}-2x+1}\right)\\
	&=&\lim\limits_{x\to-\infty}\left(\dfrac{3+\dfrac{1}{x}}{-\sqrt{4-\dfrac{1}{x}+\dfrac{2}{x^2}}-2+\dfrac{1}{x}}\right)\\
	&=&-\dfrac{3}{4}.
	\end{eqnarray*}}
\end{bt}

\begin{bt}%[1D9H1-3]%[ma:b7f0e471]
% nguon: Archive/2024/Update-SP-BG10-11-3/SP-3/BG11-3/data_1C5/DS11-C5-B2-Dang7-TinhXacSuatBangKetHopCacQuyTac-TranTony.tex #1547
%[Bài giảng 11, Trần Tony]
	Một đề trắc nghiệm môn toán có $50$ câu hỏi, mỗi câu hỏi có $4$ phương án chọn, trong đó có $1$ phương án đúng, chọn phương án đúng thì câu đó được $0{,2}$ điểm. Trong thời gian cho phép $90$ phút bạn Lân đã làm bài chắc chắn đúng $40$ câu, $10$ còn lại bạn trả lời ngẫu nhiên. Tính xác suất $p$ để bạn Lân được đúng $9$ điểm.
	\loigiai{
		Để được đúng $9$ điểm Lân phải trả lời đúng $45$ câu, trong đó có $5$ câu trả lời ngẫu nhiên, xác suất để trả lời đúng một câu là $0{,}25$. Xác suất để Lân trả lời đúng $5$ câu trong $10$ câu khi trả lời ngẫu nhiên là $p=\mathrm{C}_{10}^5  0{,}25^5 \cdot 0{,}75^5.$
	}
\end{bt}

\begin{bt}%[1D6H4-3]%[ma:2f7ccb42]
% nguon: Archive/SP-4Phan-Dot3/Data/CD_1_OTHK2_De04.tex #4302

	Giải bất phương trình $\log_4 (x+7)>\log_2 (x+1)$.
	\loigiai{Điều kiện $ x>-1 $.\\
		Ta có
		\allowdisplaybreaks
		\begin{eqnarray*}
			& & \log_4 (x+7)>\log_2 (x+1)\\
			&\Leftrightarrow &\dfrac{1}{2}\log_2 (x+7)>\log_2 (x+1)\\
			&\Leftrightarrow & \log_2 (x+7)>2\log_2 (x+1)\\
			& \Leftrightarrow &x+7>(x+1)^2\\
			&\Leftrightarrow & x^2+x-6<0\Leftrightarrow -3<x<2.
		\end{eqnarray*}
		Vậy tập nghiệm của bất phương trình $S=(-1 ; 2)$.
	}
\end{bt}

\begin{bt}%[1D2C3-8]%[ma:0b7b28f3]
% nguon: Archive/2025-Sangtac/Khoi11/Data/1-CD/CD_1-OTHK1_De10.tex #4709

	Một quả bóng được thả thẳng đứng từ độ cao $8$ mét rơi xuống đất và nảy lên. Giả sử sau mỗi một lần rơi xuống, nó nảy lên được một độ cao bằng $60\%$ độ cao vừa rơi xuống.
	Tính tổng quãng đường quả bóng di chuyển được kể từ lúc thả xuống đến khi quả bóng chạm đất lần thứ 8 (làm tròn kết quả đến hàng phần mười của mét).\par
	\shortans[0]{$19{,}7$}
	\loigiai{Gọi $u_n$ (m) là độ cao mà quả bóng đạt được sau khi nảy lên ở lần thứ $n$.\\
		Khi đó, dãy $\left(u_n\right)$ lập thành cấp số nhân có $u_1 = 8 \cdot 0{,}6 = 4{,}8$ và công bội $q = 0{,}6$. \\
		Kể từ lúc thả xuống đến khi quả bóng chạm đất lần thứ $8$, quả bóng đã được nảy lên $7$ lần rồi lại rơi xuống. Do quãng đường quả bóng nảy lên và rơi xuống bằng nhau nên tổng quãng đường quả bóng di chuyển được kể từ lúc thả xuống đến khi quả bóng chạm đất lần thứ $8$ là
		$$ S= 8 + 2 \left(u_1 + u_2 + \ldots + u_{7}\right) = 8 + 2 \cdot 4{,}8 \cdot \dfrac{1- 0{,}6^{7}}{1- 0{,}6} \approx 19{,}7 \,(\mathrm{m}).$$
	}
\end{bt}

\begin{bt}%[1D1V4-8]%[ma:7ad8903e]
% nguon: Archive/2025-Sangtac/Khoi11/Data/1-CTST/CTST_1-OTGK1_De08.tex #4869

Một quả đạn pháo được bắn ra khỏi nòng pháo với vận tốc ban đầu $ v_0=500$ m/s hợp với phương ngang một góc $ \alpha $. Trong Vật lí, ta biết rằng, nếu bỏ qua sức cản của không khí và coi quả đạn pháo được bắn ra từ mặt đất thì quỹ đạo của quả đạn tuân theo phương trình $ y=\dfrac{-g}{2v_0^2\cos^2\alpha}\cdot x^2+x\tan \alpha $, trong đó $g = 10$ m/s$ ^2 $ là gia tốc trọng trường. Góc bắn $ \alpha $ để quả đạn bay xa nhất là bao nhiêu độ?
\shortans[0]{45}
\loigiai{
Ta có \begin{eqnarray*}
y=0&\Leftrightarrow& \dfrac{-g}{2v_0^2\cos^2\alpha}\cdot x^2+x\tan \alpha=0\\
&\Leftrightarrow & x\left(\dfrac{-g}{2v_0^2\cos^2\alpha}\cdot x+\tan \alpha \right)=0 \\
&\Leftrightarrow & \hoac{&x=0\\&x=\dfrac{v_0^2\sin2\alpha}{g}= 25~000\sin2\alpha.}
\end{eqnarray*}
Vậy tầm bay xa mà quả đạn đạt tới là $25~000\sin2\alpha$ m.\\
Vậy đạn bay xa nhất khi $ \sin2\alpha=1\Leftrightarrow 2\alpha=\dfrac{\pi}{2}+k2\pi\Leftrightarrow \alpha =\dfrac{\pi}{4}+k\pi $.\\
Vì $ 0\le\alpha\le \dfrac{\pi}{2} $ nên $ \alpha=\dfrac{\pi}{4} $.\\
Vậy góc bắn bằng $ 45^\circ $ thì đạn bay xa nhất.
}
\end{bt}

\begin{bt}%[1D9V1-3]%[ma:343c2e58]
% nguon: Archive/SP_KhoLatex/SP_Đề Theo Bài-Khối 11/1-B28.tex #4146

	Hai xạ thủ cùng bắn mỗi người một viên đạn vào bia. Xác suất bắn trúng bia của hai xạ thủ lần lượt là $\dfrac{1}{3}$ và $\dfrac{1}{4}$. Biết rằng việc bắn súng của hai xạ thủ là độc lập với nhau.	Tính xác suất của biến cố cả hai xạ thủ đều bắn trúng bia? (Kết quả làm tròn đến hàng phần trăm)
	\shortans{$0{,}08$}
	\loigiai{
		Gọi $A,~B$ lần lượt là các biến cố “xạ thủ thứ nhất bắn trúng bia” và “xạ thủ thứ hai bắn trúng bia”. Khi đó $A,~B$ là hai biến cố độc lập. Áp dụng quy tắc nhân cho hai biến cố độc lập ta được:\\
		Xác suất để cả hai xạ thủ đều bắn trúng bia là:\\
		$P(AB)=P(A)\cdot P(B)=\dfrac{1}{3}\cdot \dfrac{1}{4}=\dfrac{1}{12}\approx 0{,}08$.
		
	}
\end{bt}

\begin{bt}%[1D2V3-8]%[ma:2fd9bbb8]
% nguon: Archive/SP_SauLuoi_2024-2025/SP_HKI/1-TN-DS-TLN-THPT-ChuyenHungVuong-PhuTho-HKI-NH24-25.tex #22
%[Dự án đề kiểm tra Toán khối 11 học kì 1 NH24-25-Dot 5-Thái Bảo ]%[THPT Chuyên Hùng Vương - Phú Thọ]
	Theo thống kê của Chi cục Dân số Hà Nội, tính đến năm $2024$, dân số thủ đô Hà Nội ước tính đạt khoảng $8{,}5$ triệu người và tốc độ tăng trưởng dân số là $1{,}26\%$. Nếu tốc độ tăng trưởng dân số này được giữ nguyên hàng năm, hãy ước tính dân số của thủ đô Hà Nội vào năm $2030$ (tính theo đơn vị triệu người, làm tròn đến hàng phần trăm).
	\shortans{$9{,}16$}
	\loigiai{
		Theo giả thuyết đề bài, dân số thủ đô Hà Nội qua các năm là một cấp số nhân với $u_1=8{,}5$ và $q=1{,}26\%$.\\
		Dân số $2030$ là số hạng thứ $7$ của cấp số nhân trên. Do đó
		$$u_7=u_1\cdot q^6=8{,}5\cdot \left(1+1{,}26\%\right)^6\approx 9{,}16.$$
	}
\end{bt}

\begin{ex}%[1D2H2-7]%[ma:b845add3]
% nguon: Archive/SP_SauLuoi_2024-2025/SP_GHKI/1-TL-TN-DS-NguyenThiMinhKhai-HCM-GHKI-NH24-25.tex #126
%[Dự án đề kiểm tra Toán 11 GHKI NH24-25-Dot3-Hieu Hieu Minh Minh]%[THPT Nguyễn Thị Minh Khai - Tp HCM]
	Bạn Lam muốn mua một máy tính Casio Fx-580VNX tại một nhà sách có giá là $840$ ngàn đồng. Bạn Lam lập kế hoạch mỗi ngày sẽ tiết kiệm $20$ ngàn đồng tiền tiêu vặt mà ba mẹ cho. Gọi $u_n$ ( $n \in \mathbb{N}^*$ ) là số tiền bạn Lam có được vào ngày thứ $n$, biết rằng ngày đầu tiên bạn Lam có được số tiền là $70$ ngàn đồng.
	\choiceTF[t]
	{\True Dãy số $\left(u_n\right)$ là một cấp số cộng có số hạng đầu $u_1=70$ (ngàn đồng) và công sai $d=20$ (ngàn đồng)}
	{Công thức số hạng tổng quát của dãy số $\left(u_n\right)$ là $u_n=90-20 n$ (ngàn đồng), $n \in \mathbb{N}^*$}
	{Bạn Lam có được số tiền là $270$ (ngàn đồng) vào ngày thứ $10$}
	{\True Vào ngày thứ $40$, bạn Lam đã đủ tiền mua máy tính Casio Fx580 VNX tại nhà sách trên}
	\loigiai{
		
		Ta có $u_1=70$ và công sai $d=20$ nên $u_n=70+(n-1)\cdot 20=20n+50$ (ngàn đồng).
		\begin{itemchoice}
			\itemch \textbf{Đúng}. Theo chứng minh trên thì $\left(u_n\right)$ là một cấp số cộng với số hạng đầu $u_1=70$ (ngàn đồng) và công sai $d=20$ (ngàn đồng)  
			\itemch \textbf{Sai}. Theo chứng minh trên thì $u_n=70+(n-1)\cdot 20=20n+50$ (ngàn đồng).
			\itemch \textbf{Sai}. Ngày thứ $10$, bạn Lam có số tiền là $u_{10}=20\cdot 10+50=250$ ngàn đồng.
			\itemch \textbf{Đúng}. \\
			Ngày thứ $40$, bạn Lam có số tiền là $u_{40}=20\cdot 40+50=850$ ngàn đồng.\\
			Vì $850$ ngàn đồng > $840$ ngàn đồng nên bạn Lam đã đủ tiền mua máy tính Casio Fx580 VNX tại nhà sách trên.
		\end{itemchoice}
	}
\end{ex}

\begin{ex}%[1D5N1-2]%[ma:3a65af31]
% nguon: Archive/2024/Update-SP-BG10-11-3/SP-3/BG11-3/data_1C5/ChuongV-Bai1-Dang1-MauSoLieuGhepNhom-Diamon.tex #1532
%[Dự án EX-TF-2024, Mui Doan]
Điều tra về chiều cao của $100$ học sinh lớp $11$ của một trường THPT, ta được kết quả
\begin{center}
\begin{tabular}{|c|c|c|c|c|c|c|c|}
\hline \begin{tabular}{c} 
Chiều cao \\
(cm)
\end{tabular} & {$[150 ; 152)$} & {$[152 ; 154)$} & {$[154 ; 156)$} & {$[156 ; 158)$} & {$[158 ; 160)$} & {$[160 ; 162)$} & {$[162 ; 168)$} \\
\hline Số học sinh & $5$ & $18$ & $40$ & $25$ & $8$ & $3$ & $1$ \\
\hline
\end{tabular}
\end{center}
Các khẳng định sau đúng hay sai?
\choiceTF[t]
{\True Mẫu số liệu trên có $7$ nhóm}
{Số học sinh có chiều cao dưới $154$ cm là $18$}
{\True Số học sinh có chiều cao trên $160$ cm là $4$}
{\True Số học sinh có chiều cao trong khoảng từ $156$ cm đến dưới $160$ cm là $33$}
\loigiai{
\begin{itemchoice}
\itemch  Mẫu số liệu trên có $7$ nhóm.
\itemch  Số học sinh có chiều cao dưới $154$ cm là $5+18=23$.
\itemch   Số học sinh có chiều cao trên $160$ cm là $3+1=4$.
\itemch  Số học sinh có chiều cao trong khoảng từ $156$ cm đến dưới $160$ cm là $25+8=33$.
\end{itemchoice}
}
\end{ex}

\begin{ex}%[1D1H1-3]%[ma:8d843e85]
% nguon: Archive/2024/TOAN_11_XE/TOAN_11_XE/data4/11C1/BTTN1.tex #1064
%[Trương Đăng Khoa - TeX Đúng Sai Toán 11]
	Xét tính đúng - sai của các mệnh đề sau.
	\choiceTF
	{\True Cho góc hình học $\widehat{u O v}=75^{\circ}$. Công thức số đo tổng quát của góc lượng giác $(O u, O v)$ là $75^{\circ}+k 360^{\circ}, k \in \mathbb{Z}$}
	{Cho một góc lượng giác $(O x, O u)$ có số đo $50^{\circ}$ và một góc lượng giác $(O x, O v)$ có số đo $-125^{\circ}$. Số đo của góc lượng giác $(O u, O v)$ bằng $175^{\circ}+k 360^{\circ}$, $k \in \mathbb{Z}$}
	{Trên đồng hồ hình dưới đây, kim phút đang chỉ đúng số $2$. Đến lúc $13$ giờ $30$ phút thì kim phút đã quay được một góc lượng giác có số đo bằng $480^{\circ}$\\
		\centerline{\begin{tikzpicture}[scale=1, font=\footnotesize, line join=round, line cap=round,>=stealth]
				\coordinate (O) at (0,0);
				\draw (O) circle (3 cm);
				%\fill[gray!50] (O) circle ( 1 cm);
				%Trục Oxy.
				\coordinate (3) at ($(O) + (0:2.5)$);
				\coordinate (2) at ($(O) + (30:2.5)$);
				\coordinate (1) at ($(O) + (60:2.5)$);
				\coordinate (12) at ($(O) + (90:2.5)$);
				\coordinate (11) at ($(O) + (120:2.5)$);
				\coordinate (10) at ($(O) + (150:2.5)$);
				\coordinate (9) at ($(O) + (180:2.5)$);
				\coordinate (8) at ($(O) + (210:2.5)$);
				\coordinate (7) at ($(O) + (240:2.5)$);
				\coordinate (6) at ($(O) + (270:2.5)$);
				\coordinate (5) at ($(O) + (300:2.5)$);
				\coordinate (4) at ($(O) + (330:2.5)$);
				\draw[line width=3pt, ->] (O)--(0.1,1.3);
				\draw[line width=3pt, ->]  (O)--(2);
				\foreach \i/\g in {3/0, 2/60,1/60,12/90,11/120,10/120,9/180,8/-120,7/-120,6/-90,5/-45,4/-45}{\draw[fill=black](\i) circle (2pt) ($(\i)+(\g:3mm)$) node[scale=1]{$\i$};}
				\draw (0,-3.7) node {$12$ giờ $30$ phút};
		\end{tikzpicture}}
	}
	{\True Trong hình dưới, cánh quạt được chia thành $3$ phần bằng nhau  $\widehat{A O D}=28^{\circ}$.  Công thức số đo tồng quát của góc lượng giác $(O x, O C)$ bằng $-148^{\circ}+k \cdot 360^{\circ}$\\
		\centerline{\begin{tikzpicture}[scale=1, font=\footnotesize, line join=round, line cap=round,>=stealth]
				\coordinate (O) at (0,0);
				\draw[dashed] (O) circle (2.5 cm);
				\fill[gray!50] (O) circle ( 1 cm);
				%Trục Oxy.
				\draw[->] (-3.5,0)--(3.5,0) node[below]{$x$};
				\draw[->] (0,-3.5)--(0,3.5) node[left]{$y$};
				\coordinate (A) at ($(O) + (0:2.5)$);
				\coordinate (D) at ($(O) + (-30:2.5)$);
				\coordinate (C) at ($(O) + (210:2.5)$);
				\coordinate (B) at ($(O) + (90:2.5)$);
				\draw[line width=3pt] (O)--(C) (O)--(B) (O)--(D);
				\foreach \i/\g in {A/45,B/45,C/-120}{\draw[fill=white](\i) circle (1.5pt) ($(\i)+(\g:3mm)$) node[scale=1]{$\i$};}
		\end{tikzpicture}}
	}
	\loigiai{
		\begin{itemchoice}
			\itemch Đúng.
			\itemch  Sai. $ \text{sđ}( {Ou},  {Ov})= \text{sđ}( {Ox},  {Ov})- {s} d( {Ox},  {Ou})+ {k} 360^{\circ}=-175^{\circ}+k 360^{\circ}$, $k \in \mathbb{Z}$.
			\itemch  Sai. Kim phút đã quay được một góc lượng giác có số đo bằng \break  $-\left(360^{\circ}+120^{\circ}\right)=-480^{\circ}$.
			\itemch  Đúng. Cánh quạt được chia thành $3$ phần bằng nhau, vậy mỗi phần sẽ có số đo là $\dfrac{360^{\circ}}{ 3}=120^{\circ}$.\\
			Số đo góc lượng giác $(O D, O C)=-120^{\circ}$.\\
			Theo hệ thức Chasles ta có
			$$
			(O x, O C)=(O x, O D)+(O D, O C)=-28^{\circ}+\left(-120^{\circ}\right)+k 360^{\circ}=-148^{\circ}+k 360^{\circ} .
			$$
		\end{itemchoice}
	}
\end{ex}